\documentclass[11pt]{article}
\usepackage[margin=1in]{geometry}
\usepackage[T1]{fontenc}
\usepackage[utf8]{inputenc}
\usepackage{lmodern}
\usepackage{microtype}
\usepackage{setspace}
\usepackage{xcolor}
\usepackage{enumitem}
\usepackage{graphicx}
\usepackage{booktabs}
\usepackage{array}
\usepackage{float}
\usepackage{listings}
\usepackage{xurl}
\usepackage{hyperref}
\hypersetup{
  colorlinks=true,
  linkcolor=blue,
  citecolor=blue,
  urlcolor=blue,
  breaklinks=true
}
\setstretch{1.08}
\urlstyle{same}
\setlength{\emergencystretch}{2em}

% Listings style for log output / code snippets
\lstset{
  basicstyle=\ttfamily\small,
  breaklines=true,
  frame=single,
  backgroundcolor=\color{gray!8},
  captionpos=b
}

\newcommand{\todo}[1]{\textcolor{red}{[TODO: #1]}}
\newcommand{\rpc}[1]{\texttt{#1}}

\title{\textbf{Distributed Consensus in a Multiplayer Fishing Game:\\
Two-Phase Commit and Raft over a 6-Node gRPC Cluster}}
\author{
  Joe Caldwell\\
  \small Student ID: \todo{insert student ID}
  \and
  \todo{Team Member 2 Name}\\
  \small Student ID: \todo{insert student ID}
  % Add more \and blocks if needed
}
\date{\today}

\begin{document}
\maketitle

%--------------------------------------------------------------------
\section*{Project Link}
%--------------------------------------------------------------------
\textbf{Source code (GitHub):}
\url{https://github.com/Mojo4Sho1/CSCE5306_Project3}

%--------------------------------------------------------------------
\section*{Work Division}
%--------------------------------------------------------------------
\begin{tabular}{p{0.35\linewidth} p{0.60\linewidth}}
\toprule
\textbf{Team Member} & \textbf{Contributions} \\
\midrule
Joe Caldwell & \todo{describe contributions} \\
\todo{Member 2} & \todo{describe contributions} \\
\bottomrule
\end{tabular}

%--------------------------------------------------------------------
\section{Introduction}
%--------------------------------------------------------------------
This report describes the extension of a baseline multiplayer fishing game with two distributed
consensus algorithms: simplified Two-Phase Commit (2PC) and the Raft consensus algorithm.
The baseline system is a Python gRPC prototype that runs six containerized nodes
and allows players to log in, update their location, fish, and view other players.

The project adds replicated state management to the \rpc{UpdateLocation} RPC---the primary
mutation in the game---via two complementary mechanisms:
\begin{enumerate}[label=(\alph*)]
  \item \textbf{Two-Phase Commit (Q1 + Q2):} a coordinator node drives a voting phase
    followed by a global-commit or global-abort decision.  All nodes participate as
    2PC participants.  Inter-node communication uses gRPC; the coordinator also issues
    intra-node gRPC calls to connect the voting and decision phases within the same container.
  \item \textbf{Raft (Q3 + Q4):} a leader-election state machine selects a single leader.
    The leader replicates location-update operations to followers via the log;
    a majority acknowledgement is required before an entry is committed and applied to state.
    Non-leader nodes forward client requests to the current leader via a \rpc{ForwardRequest} RPC.
\end{enumerate}

Both mechanisms apply to the same \rpc{UpdateLocation} RPC, making it straightforward to
observe their effects through the game's existing client.
Fault-tolerance behaviour is evaluated through five failure scenarios (Q5) that exercise
leader crash, follower crash, network partition, new-node join, and split-vote re-election.

%--------------------------------------------------------------------
\section{System Architecture}
%--------------------------------------------------------------------

\subsection{Baseline}

The baseline is a six-node Python gRPC cluster defined in \texttt{server/docker-compose.yml}.
Each node (\texttt{fishing1}--\texttt{fishing6}) runs the same \texttt{server.py} on ports
50051--50056.  The client-facing API is defined in \texttt{fishing.proto} (not modified) and
exposes seven RPCs: \rpc{Login}, \rpc{UpdateLocation} (client-streaming),
\rpc{ListUsers}, \rpc{StartFishing}, \rpc{CurrentUsers}, \rpc{Inventory}, and \rpc{GetImage}.
State (users and inventory) is in-memory.

\subsection{Extension Overview}

The extension adds two proto files and supporting Python modules without modifying
\texttt{fishing.proto} or the existing generated stubs.
Table~\ref{tab:components} lists the new components.

\renewcommand{\arraystretch}{1.12}
\begin{table}[H]
\centering
\caption{New components added by this project.}
\label{tab:components}
\begin{tabular}{p{0.22\linewidth} p{0.68\linewidth}}
\toprule
\textbf{File / Module} & \textbf{Purpose} \\
\midrule
\texttt{server/twopc.proto}          & 2PC service definitions: \rpc{VoteRequest}, \rpc{VoteResponse}, \rpc{GlobalDecision}, \rpc{DecisionAck}, intra-node \rpc{ReportVote} / \rpc{NotifyDecision} \\
\texttt{server/raft.proto}           & Raft service definitions: \rpc{RequestVote}, \rpc{AppendEntries}, \rpc{ForwardRequest} and associated messages including \texttt{LogEntry} \\
\texttt{server/raft\_node.py}        & \texttt{RaftNode} state machine (follower / candidate / leader), \texttt{RaftServicer} gRPC handler, log replication helpers \\
\texttt{server/server.py}            & Extended with 2PC coordinator/participant servicers, Raft wiring, and \rpc{UpdateLocation} routing \\
\texttt{tests/unit/}                 & 55 unit tests covering 2PC, Raft election, and Raft log replication (no Docker required) \\
\bottomrule
\end{tabular}
\end{table}
\renewcommand{\arraystretch}{1.0}

\subsection{Node Configuration}

All six nodes run the same server image.
Each node is identified by the \texttt{NODE\_ID} environment variable (1--6) and discovers
its peers via the \texttt{PEERS} variable (comma-separated \texttt{host:port} strings).
One node (\texttt{fishing1}) is designated the 2PC coordinator via \texttt{IS\_COORDINATOR=true}.
The Raft leader is elected dynamically; any node may become the leader.

%--------------------------------------------------------------------
\section{Two-Phase Commit (Q1 + Q2)}
%--------------------------------------------------------------------

\subsection{Protocol Design}

When the coordinator node (\texttt{fishing1}) receives an \rpc{UpdateLocation} RPC it
drives a two-phase commit over all other nodes:

\begin{enumerate}
  \item \textbf{Voting phase}: the coordinator sends a \rpc{VoteRequest} to each participant.
    Participants vote \emph{commit} (coordinates are non-negative) or \emph{abort}
    (negative $x$ or $y$).  After collecting all votes, the coordinator calls the
    intra-node \rpc{ReportVote} RPC on its own decision-phase servicer (localhost) to
    hand off the aggregated result.
  \item \textbf{Decision phase}: based on unanimous commit or any abort, the coordinator
    sends a \rpc{GlobalDecision} to each participant.  Participants apply
    \texttt{state.update\_user} on global-commit or discard on global-abort.
    The decision phase then calls intra-node \rpc{NotifyDecision} to inform the voting
    phase of the final outcome.
\end{enumerate}

The intra-node gRPC calls (coordinator $\rightarrow$ localhost) are required by the
assignment to demonstrate that the voting and decision phases are distinct services
communicating via gRPC even within the same container.

\subsection{RPC Log Format}

All 2PC RPC calls print the assignment-required log lines.
Table~\ref{tab:2pc-log} shows examples for a commit scenario on a 3-node excerpt.

\begin{table}[H]
\centering
\caption{Representative 2PC log lines (voting + decision phases, node 1 as coordinator).}
\label{tab:2pc-log}
\begin{tabular}{p{0.92\linewidth}}
\toprule
\textbf{Log line} \\
\midrule
\lstinline|Phase voting of Node 1 sends RPC VoteRequest to Phase voting of Node 2| \\
\lstinline|Phase voting of Node 2 sends RPC VoteRequest to Phase voting of Node 1| \\
\lstinline|Phase voting of Node 1 sends RPC ReportVote to Phase decision of Node 1| \\
\lstinline|Phase decision of Node 1 sends RPC GlobalDecision to Phase decision of Node 2| \\
\lstinline|Phase decision of Node 2 sends RPC GlobalDecision to Phase decision of Node 1| \\
\lstinline|Phase decision of Node 1 sends RPC NotifyDecision to Phase voting of Node 1| \\
\bottomrule
\end{tabular}
\end{table}

\subsection{Abort Semantics}

A participant votes \emph{abort} when the proposed location update has a negative $x$
or $y$ coordinate.  If any participant votes abort the coordinator broadcasts
\rpc{GlobalDecision(global\_commit=false)} and no node applies the update.

%--------------------------------------------------------------------
\section{Raft Leader Election (Q3)}
%--------------------------------------------------------------------

\subsection{State Machine}

Each node starts as a \emph{follower} with \texttt{current\_term = 0}.
An election timeout in the range $[1.5\text{s}, 3.0\text{s}]$ (randomised per node)
drives follower $\rightarrow$ candidate transitions.
A candidate that wins a majority becomes the \emph{leader} for its term and sends
periodic \rpc{AppendEntries} heartbeats every 1 second.

\begin{itemize}
  \item \textbf{RequestVote}: a candidate sends \rpc{RequestVote} to all peers.
    A node grants its vote if it has not yet voted in the current term (or if it
    previously voted for the same candidate) and the request term is not stale.
    A higher-term request causes any node (including a leader) to step down to follower.
  \item \textbf{Majority}: a candidate becomes leader if it collects votes from
    $> N/2$ nodes (including its self-vote), where $N$ is the total node count.
    For the 6-node cluster, 4 votes are required.
\end{itemize}

\subsection{RPC Log Format}

\begin{table}[H]
\centering
\caption{Raft RPC log format (election phase).}
\label{tab:raft-log-election}
\begin{tabular}{p{0.92\linewidth}}
\toprule
\textbf{Log line} \\
\midrule
\lstinline|Node 1 sends RPC RequestVote to Node 2| \\
\lstinline|Node 2 runs RPC RequestVote called by Node 1| \\
\lstinline|Node 1 sends RPC AppendEntries to Node 3| \\
\lstinline|Node 3 runs RPC AppendEntries called by Node 1| \\
\bottomrule
\end{tabular}
\end{table}

%--------------------------------------------------------------------
\section{Raft Log Replication (Q4)}
%--------------------------------------------------------------------

\subsection{Log Entry Structure}

Each log entry is a triple $\langle \textit{operation}, \textit{term}, \textit{index} \rangle$:
\begin{itemize}
  \item \textit{operation}: a string encoding of the update command,
    e.g.\ \lstinline|UpdateLocation:alice:pass:10.5:20.3|.
  \item \textit{term}: the leader's current term when the entry was created.
  \item \textit{index}: a 1-based sequential integer.
\end{itemize}

\subsection{Replication Protocol}

\begin{enumerate}
  \item \textbf{Leader receives client request}: if the node is the leader, it calls
    \texttt{append\_log\_entry(operation)}, which appends the entry, records its own
    acknowledgement in an \texttt{ack\_tracker} set, and creates a \texttt{threading.Event}
    for the waiting client.
  \item \textbf{Heartbeat with full log}: on every heartbeat (every 1 second) the leader
    sends its \emph{entire} log plus the current \texttt{commit\_index} in the
    \rpc{AppendEntries} request.  The simplified Raft variant used here sends the full
    log rather than incremental entries.
  \item \textbf{Follower receives AppendEntries}: the follower replaces its log with the
    received entries, then applies all entries with
    $\textit{index} \le \textit{commit\_index}$ that have not yet been applied
    (\texttt{last\_applied}).  It returns \texttt{success=true} to ACK.
  \item \textbf{Leader commits}: when the leader receives ACKs from a majority of nodes
    (including itself) for a given log index, it commits that entry, applies the operation
    to local state, increments \texttt{commit\_index}, and sets the client's event to
    unblock the waiting \rpc{UpdateLocation} call.  A 5-second timeout prevents clients
    from hanging indefinitely if the leader steps down before committing.
  \item \textbf{Non-leader forwarding}: if a client connects to a follower, the follower's
    \rpc{UpdateLocation} handler routes the request to the leader via a
    \rpc{ForwardRequest} RPC.
    The leader's \rpc{ForwardRequest} handler appends the entry and waits for commit
    before responding.
\end{enumerate}

\subsection{ACK Deduplication}

ACKs are tracked per-entry using a \texttt{set} of follower IDs
(\texttt{ack\_tracker[\textit{index}] = \{follower\_ids\}}).
Because the leader sends the full log on every heartbeat, the same follower may send
multiple ACKs for the same entry across heartbeat rounds.
The set-based tracker prevents double-counting.

\subsection{RPC Log Format}

\begin{table}[H]
\centering
\caption{Raft RPC log format (log replication phase).}
\label{tab:raft-log-replication}
\begin{tabular}{p{0.92\linewidth}}
\toprule
\textbf{Log line} \\
\midrule
\lstinline|Node 1 sends RPC AppendEntries to Node 4| \\
\lstinline|Node 4 runs RPC AppendEntries called by Node 1| \\
\lstinline|Node 3 sends RPC ForwardRequest to Node 1| \\
\lstinline|Node 1 runs RPC ForwardRequest called by Node 3| \\
\bottomrule
\end{tabular}
\end{table}

%--------------------------------------------------------------------
\section{Failure Tests (Q5)}
%--------------------------------------------------------------------

All five test cases are run against the live 6-node Docker cluster started with
\texttt{make up}.  The cluster is reset (\texttt{make down \&\& make up}) between test cases.

%-------------------------------------------------------------------
\subsection{Test Case 1: Leader Crash and Re-Election}
%-------------------------------------------------------------------

\paragraph{Setup.}
Start the cluster and wait for an initial leader to be elected.
Send several \rpc{UpdateLocation} requests to confirm normal operation.

\paragraph{Action.}
\begin{lstlisting}[caption={TC1: kill the leader container.}]
docker stop fishing<leader_id>
\end{lstlisting}

\paragraph{Expected behaviour.}
Remaining followers detect a missing heartbeat after their randomised election timeout
(1.5--3 s).  One follower transitions to candidate, issues \rpc{RequestVote} RPCs, wins
a majority, and begins sending heartbeats as the new leader.
The client can resume \rpc{UpdateLocation} requests via the new leader (or via forwarding
from any remaining follower).

\paragraph{Observed behaviour.}
\todo{Fill in after running the test. Note which node was killed, how long re-election took, and whether the client resumed successfully.}

\begin{figure}[H]
  \centering
  \includegraphics[width=\linewidth]{screenshots/tc1_leader_crash.png}
  \caption{TC1 --- terminal output showing leader crash, election process, and resumed operations.}
  \label{fig:tc1}
\end{figure}

%-------------------------------------------------------------------
\subsection{Test Case 2: Follower Crash and Recovery}
%-------------------------------------------------------------------

\paragraph{Setup.}
Start the cluster, elect a leader, send several updates to populate the log.

\paragraph{Action.}
\begin{lstlisting}[caption={TC2: kill a follower, send updates, restart follower.}]
docker stop fishing<follower_id>
# send more UpdateLocation requests from the client
docker start fishing<follower_id>
\end{lstlisting}

\paragraph{Expected behaviour.}
The cluster continues with five active nodes; the leader commits entries with ACKs from
the remaining four (majority of six).  When the follower restarts it receives the full
leader log on the next heartbeat and catches up to the current \texttt{commit\_index}.

\paragraph{Observed behaviour.}
\todo{Fill in after running the test. Note how many updates were sent while the follower was down and confirm the restarted follower's log matched the leader's log.}

\begin{figure}[H]
  \centering
  \includegraphics[width=\linewidth]{screenshots/tc2_follower_recovery.png}
  \caption{TC2 --- logs showing continued operation during follower absence and log catch-up after restart.}
  \label{fig:tc2}
\end{figure}

%-------------------------------------------------------------------
\subsection{Test Case 3: Network Partition (Pause / Unpause)}
%-------------------------------------------------------------------

\paragraph{Setup.}
Start the cluster with an established leader and several log entries.

\paragraph{Action.}
\begin{lstlisting}[caption={TC3: pause a container to simulate a network partition.}]
docker pause fishing<node_id>
# wait 5-10 seconds; send UpdateLocation requests
docker unpause fishing<node_id>
\end{lstlisting}

\paragraph{Expected behaviour.}
If the paused node is the leader, the remaining nodes detect a missing heartbeat and
elect a new leader.  If it is a follower, the cluster continues normally.
On unpause the node receives an \rpc{AppendEntries} heartbeat, discovers the current term,
and syncs its log.  A stale leader that unpauses steps down upon seeing a higher term.

\paragraph{Observed behaviour.}
\todo{Fill in after running the test. State whether the leader or a follower was paused and describe the observed resynchronisation.}

\begin{figure}[H]
  \centering
  \includegraphics[width=\linewidth]{screenshots/tc3_partition.png}
  \caption{TC3 --- terminal output showing the pause, cluster reaction, unpause, and resynchronisation.}
  \label{fig:tc3}
\end{figure}

%-------------------------------------------------------------------
\subsection{Test Case 4: New Node Joining the Cluster}
%-------------------------------------------------------------------

\paragraph{Setup.}
Start only five nodes (\texttt{fishing1}--\texttt{fishing5}), wait for leader election,
and send several updates so the log contains committed entries.

\paragraph{Action.}
\begin{lstlisting}[caption={TC4: start the sixth node after the cluster is already running.}]
docker compose -f server/docker-compose.yml up -d fishing6
\end{lstlisting}

\paragraph{Expected behaviour.}
\texttt{fishing6} starts as a follower.
On the next heartbeat from the leader it receives the full log via \rpc{AppendEntries},
applies all committed entries, and joins the cluster with its state consistent with the
other nodes.

\paragraph{Observed behaviour.}
\todo{Fill in after running the test. Confirm that fishing6 received AppendEntries and that its log matched the leader's log.}

\begin{figure}[H]
  \centering
  \includegraphics[width=\linewidth]{screenshots/tc4_new_node.png}
  \caption{TC4 --- logs showing the new node startup, log sync via AppendEntries, and cluster membership.}
  \label{fig:tc4}
\end{figure}

%-------------------------------------------------------------------
\subsection{Test Case 5: Split Vote and Election Retry}
%-------------------------------------------------------------------

\paragraph{Setup.}
Start the cluster and identify the current leader.

\paragraph{Action.}
Kill the leader and one follower simultaneously to increase the probability of a split vote
among the remaining four nodes.
\begin{lstlisting}[caption={TC5: kill leader and one follower simultaneously.}]
docker stop fishing<leader_id> fishing<follower_id>
\end{lstlisting}

\paragraph{Expected behaviour.}
Multiple candidates may start elections in the same term.
If no candidate receives a majority (3 of the remaining 4 nodes), the term ends without
a leader.  Each candidate re-randomises its election timeout and starts a new election
with an incremented term.  Eventually one candidate wins a majority and becomes leader.

\paragraph{Observed behaviour.}
\todo{Fill in after running the test. If a genuine split vote was not observed, describe the election timeout mechanism from the logs (different timeouts per node) and show that multiple terms were observed before a leader was elected.}

\begin{figure}[H]
  \centering
  \includegraphics[width=\linewidth]{screenshots/tc5_split_vote.png}
  \caption{TC5 --- logs showing multiple \rpc{RequestVote} rounds, term increments, and final leader selection.}
  \label{fig:tc5}
\end{figure}

%--------------------------------------------------------------------
\section{Build and Run Instructions}
%--------------------------------------------------------------------

\subsection{Prerequisites}

\begin{itemize}
  \item Python 3.11+, \texttt{grpcio}, \texttt{grpcio-tools} (installed via \texttt{make install})
  \item Docker Desktop (Linux containers)
  \item GNU Make
\end{itemize}

\subsection{Quick Start}

\begin{lstlisting}[caption={Build and start the 6-node cluster.}]
# Install dev dependencies
make install

# (Optional) regenerate proto stubs
make proto

# Run unit tests (no Docker required)
make test

# Start the 6-node Docker cluster
make up

# Run the interactive client (cluster must be running)
cd client && python3 client.py

# Stream logs from all nodes
make logs

# Tear down
make down
\end{lstlisting}

\subsection{Unusual Notes}

\begin{itemize}
  \item The 2PC proto file is named \texttt{twopc.proto} (not \texttt{2pc.proto}).
    Python cannot import a module whose name starts with a digit;
    \texttt{grpc\_tools.protoc} would generate \texttt{2pc\_pb2.py} which is an invalid
    Python identifier.
  \item All containers set \texttt{PYTHONUNBUFFERED=1} to prevent Python's default
    stdout buffering from suppressing log output in \texttt{docker logs}.
  \item Unit tests pass \texttt{\_start\_threads=False} to \texttt{RaftNode} to suppress
    background election-timer and heartbeat threads; without this flag tests are flaky.
  \item The Raft implementation uses a simplified protocol that sends the entire log on
    every heartbeat rather than incremental entries; followers replace their log wholesale.
\end{itemize}

\subsection{External Sources}

\begin{itemize}
  \item Ongaro, D.\ and Ousterhout, J.\ (2014).
    \textit{In Search of an Understandable Consensus Algorithm (Extended Version)}.
    USENIX ATC.
    \url{https://raft.github.io/raft.pdf}
  \item Gray, J.\ and Lamport, L.\ (2006).
    Consensus on Transaction Commit.
    \textit{ACM Transactions on Database Systems}, 31(1):133--160.
  \item gRPC Python documentation: \url{https://grpc.io/docs/languages/python/}
  \item Protocol Buffers language guide (proto3):
    \url{https://protobuf.dev/programming-guides/proto3/}
\end{itemize}

%--------------------------------------------------------------------
\section{Analysis and Discussion}
%--------------------------------------------------------------------

\subsection{2PC vs.\ Raft: Trade-offs}

Both protocols provide replicated location updates but under fundamentally different
consistency models.

\paragraph{Two-Phase Commit.}
2PC requires unanimous agreement: a single participant's abort aborts the entire
transaction.
The coordinator is a single point of failure; if it crashes between the voting and
decision phases, participants are blocked waiting for a decision.
In the 6-node implementation the coordinator (\texttt{fishing1}) handles all client
writes, which concentrates load on one node.
On the other hand, 2PC is straightforward to reason about and its atomicity guarantees
are strong.

\paragraph{Raft.}
Raft tolerates the failure of up to $\lfloor (N-1)/2 \rfloor$ nodes ($N=6$ $\Rightarrow$
2 node failures) while continuing to commit new entries.
Leadership is distributed: any node can become leader after an election.
The cost is additional complexity: the election state machine, term management, and the
heartbeat mechanism introduce more moving parts than 2PC.
The blocking client pattern (\texttt{wait\_for\_commit} with a 5-second timeout) also
means that a leader stepping down mid-commit will stall clients briefly.

\paragraph{Coexistence in this implementation.}
The 2PC and Raft paths are both wired to \rpc{UpdateLocation} but operate independently.
The 2PC path fires only on the designated coordinator node; all other nodes use the Raft
path.  This is a simplification: a production system would choose one replication
mechanism and apply it consistently.

\subsection{Failure Mode Comparison}

\begin{table}[H]
\centering
\caption{Fault tolerance comparison between 2PC and Raft (6-node cluster).}
\label{tab:fault-comparison}
\renewcommand{\arraystretch}{1.12}
\begin{tabular}{p{0.30\linewidth} p{0.32\linewidth} p{0.32\linewidth}}
\toprule
\textbf{Failure scenario}           & \textbf{2PC}                         & \textbf{Raft} \\
\midrule
Single participant crash            & Coordinator may abort; others wait for decision & Cluster continues if majority intact \\
Coordinator crash                   & Blocked (no new transactions possible) & Leader election; new leader continues \\
Network partition (minority)        & Coordinator may be isolated; transactions block & Majority partition continues; minority stalls \\
Node rejoin after crash             & No recovery mechanism (in-memory state) & Full log replayed via heartbeat AppendEntries \\
Concurrent client requests          & Serialised through coordinator         & Leader serialises; followers forward \\
\bottomrule
\end{tabular}
\renewcommand{\arraystretch}{1.0}
\end{table}

%--------------------------------------------------------------------
\section{Lessons Learned from AI Tool Usage}
%--------------------------------------------------------------------

\todo{Fill in this section after completing the project. Describe how AI tools (ChatGPT, Claude Code, etc.) were used, what benefits were observed, and any limitations or verification steps required. The section from Project 2 provides a good template: interactive pair programming for early design, agent-driven workflow for implementation, diff-based review, and test/smoke validation as safeguards.}

%--------------------------------------------------------------------
\section{Conclusion}
%--------------------------------------------------------------------

This project extended a baseline multiplayer fishing game with two distributed consensus
mechanisms over a 6-node gRPC/Docker cluster.
The Two-Phase Commit implementation (Q1 + Q2) demonstrates atomic, coordinator-driven
replication with intra-node gRPC between phases.
The Raft implementation (Q3 + Q4) demonstrates fault-tolerant leader election, log
replication with majority commit, and client-request forwarding to the current leader.

Five failure scenarios (Q5) exercise the Raft implementation's fault tolerance:
leader crash and re-election, follower crash and recovery, network partition,
new-node joining, and split-vote election retry.

The evaluation highlights a fundamental trade-off: 2PC offers strong atomicity with simple
coordinator-driven logic but is vulnerable to coordinator failure and provides no
built-in recovery from partial failures.
Raft tolerates node failures up to the minority threshold and provides automatic recovery
through its heartbeat and log-catch-up mechanisms, at the cost of additional implementation
complexity.

\end{document}
